\begin{textsample}{POS Dim 1 – vem\_ed\_10 – Score 6.00 – vem\_ed\_10/t044.txt}  \label{ex:f1_pos_121}
Jornada de PCIs: alimentando pesquisas “Pesquisas em PCIs” trouxe exemplos de como as pesquisas acadêmicas podem \textbf{ser} incorporadas à realização de PCIs.
Ana Carolina Freschi apresentou uma síntese de suas pesquisas no mestrado e no doutorado, com abordagens relativas aos \textbf{intercâmbios} virtuais e a avaliação dos projetos (autoavaliação, pelo professor e por pares).
Citou o autor O´Dowd e trouxe uma revisão terminológica relacionada aos \textbf{intercâmbios} virtuais: telecolaboração, Teletandem (relacionados à aprendizagem de línguas estrangeiras), global virtual teams (na área de negócios), e os planos de ensino compartilhados, com \textbf{foco} em \textbf{competência} intercultural e digital, desenvolvidos na iniciativa COIL (Collaborative Online International Learning) e conhecidos na Unidade do Ensino Superior (Cesu) do Centro Paula Souza com o nome de Projetos Colaborativos Internacionais (PCIs).
Lilian de Souza (Fatec Piracicaba) apresentou sua tese sobre ensino de línguas para fins \textbf{específicos}.
A problemática tratada no doutorado se originou dos PCIs. “Trouxe para a minha pesquisa e a minha prática as atividades \textbf{que} eu desenvolvi no PCI”.
Lidiane Murad expôs um trabalho sobre os projetos de internacionalização na Fatec São José do Rio Preto, realizados desde 2015.
Ao \textbf{longo} desses anos, a unidade contabiliza 18 edições de projetos, 170 alunos participantes e 9 universidades parceiras.
Patrícia Patrício (equipe dos PCIs) sintetizou os resultados da pesquisa de percepção realizada com professores e alunos \textbf{que} participaram de PCIs no primeiro semestre de 2021.Dos 572 pesquisados, 90 \% afirmaram \textbf{que} participar de um PCI melhora desempenho acadêmico e chances no \textbf{mercado} de trabalho.
No segundo semestre de 2021, há mais de 40 projetos e 1000 alunos envolvidos, porém, a pesquisa não \textbf{foi} concluída.
Os resultados devem \textbf{ser} \textbf{analisados} no primeiro semestre de 2022.

% matched lemmas: analisar, competência, específico, foco, intercâmbio, longo, mercado, que, ser
\end{textsample}
